\scanpage{39}

d'où

\[
\boxed{\mathrm{Hom}^{d}(M^{*},\deleted{\unsure{\varphi_N}}N\otimes_k S_k)\simeq\mathrm{Hom}^{d}(\tau_{N-d}M,\tau_N(N\otimes_k S_k)}
\]

\[
\text{Si }\left\{\begin{aligned}M^{*}&\in\mathcal{M}_{-d}\\ N&\geq0\end{aligned}\right.
\]

p. ex.

\[
\boxed{\begin{aligned}
\mathrm{Hom}^{d}(M\otimes_k S_k,N\otimes_k S_k)&\xrightarrow{\sim}\mathrm{Hom}^{d}(\tau_{N-d}M_{S_k},\tau_N N_{S_k}\\
\rotatebox{90}{$\simeq$}\qquad&\\[-3pt]
\mathrm{Hom}^{\deleted{\unsure{d}}}_k(M,N)&\qquad \text{\uuline{si $d\geq0$}},N\geq0
\end{aligned}}
\]

Hom. \unsure{hom|homogènes} de degré $d\in\mathbb{Z}$

\[
\tau_\alpha M_{S_k}\longrightarrow\tau_\beta N_{S_k}\qquad(\beta\geq\alpha)
\]

a) \deleted{$d\geq0$}. Supposons d'abord $\boxed{d\geq\beta-\alpha}$ i.e $\alpha+d\geq\beta$

On a \deleted{$\mathrm{Hom}^{i}(M$}-

\[
\begin{aligned}
\mathrm{Hom}^{d}(\tau_\alpha M_{S_k},\tau_\beta N_{S_k})&\overset{\sim}{\longleftrightarrow}\mathrm{Hom}^{d}(\tau_\alpha M_{S_k},\underbrace{\tau_{\alpha+i}\tau_\beta N_{S_k}}_{=\tau_{\alpha+d}N_{S_k}})\\
&\xleftarrow{\sim}\mathrm{Hom}^{d}(M_{S_k},N_{S_k})\simeq\mathrm{Hom}_k(M,N)
\end{aligned}
\]

(Si $u:M\mathord{-} N$, il lui correspond l'hom

\[
\underset{j\geq\alpha}{mT^{(j)}}\longmapsto u(m)T^{(d)}T^{(j)}=C_{d,j}\,u(m)T^{(d+j)}\qquad\bigg)
\]

b) Supposons $0\leq d<\beta-\alpha$ \inserted{\deleted{\unsure{d'où|i.e.}} $\beta\geq\beta-d\unsure{>|\geq}\alpha$}, on a alors

\deleted{\[
\begin{aligned}
\mathrm{Hom}^{d}(\tau_\alpha M_{S_k},\tau_\beta N_{S_k})&\underset{\unsure{\sim|\text{inj.}}}{\longrightarrow}\mathrm{Hom}^{d}(\tau_{\beta-d}M_{S_k},\tau_\beta N_{S_k})\\
&\hspace{40mm}\rotatebox{90}{$\simeq$}\\[-4pt]
&\hspace{36mm}\mathrm{Hom}_k(M,N)
\end{aligned}
\]}

% archive sheet number: 38
